\documentclass[stu,12pt,floatsintext]{apa7}
% =====================================================================
%  Documento del Proyecto Final - Farmacia con Consultorio Medico
%  Formato APA 7 (student). Las figuras son CAPTURAS DE PANTALLA del
%  trabajo realizado; este documento NO incluye codigo fuente.
% =====================================================================

\usepackage[spanish,es-noindentfirst]{babel}
\usepackage[utf8]{inputenc}
\usepackage{graphicx}
\usepackage{csquotes}

% --- Metadatos APA 7 (edita estos valores) ---
\title{Farmacia con Consultorio Médico: Diseño e Implementación de una Base de Datos Relacional}
\shorttitle{Base de Datos: Farmacia con Consultorio}
\author{[Tu Nombre]}
\affiliation{[Institución]}
\course{LSTI2311: Bases de Datos}
\professor{[Nombre del Profesor]}
\duedate{\today}

\begin{document}

\maketitle

% =====================================================================
\section{Introducción}
% =====================================================================
El presente documento describe el diseño y la implementación de una base de
datos relacional para la gestión de consultas médicas y medicamentos recetados
en un consultorio con farmacia asociada. Se documenta el modelo conceptual, el
modelo lógico, la normalización aplicada y la implementación física en MySQL,
así como las vistas de reporte requeridas.

% =====================================================================
\section{Caso de Negocio}
% =====================================================================
El consultorio no cuenta con un sistema integrado para registrar consultas,
recetas, inventario de la farmacia y ventas. La solución propuesta es una base
de datos relacional que sirva de fundamento para un sistema de gestión, con
control de historial clínico, inventario en tiempo real y reportes de consumo.

% =====================================================================
\section{Modelo Entidad-Relación}
% =====================================================================
El modelo se diseñó en dbdiagram.io y contempla 20 tablas que cubren pacientes,
médicos, empleados, historiales clínicos, condiciones (alergias y
contraindicaciones), consultas, diagnósticos, medicamentos, interacciones,
recetas, proveedores, órdenes de compra, lotes, inventario, ventas y auditoría.

\begin{figure}
    \caption{Modelo entidad-relación en dbdiagram.io}
    \centering
    % Reemplaza por la captura del diagrama ER
    % \includegraphics[width=\linewidth]{img/modelo_der.png}
    \fbox{\parbox[c][6cm][c]{\linewidth}{\centering [Captura de pantalla: diagrama ER]}}
    \label{fig:der}
\end{figure}

% =====================================================================
\section{Modelo Lógico y Normalización}
% =====================================================================
Se especificaron los tipos de datos, claves primarias y foráneas de cada tabla.
Se aplicaron las tres primeras formas normales para eliminar redundancias y
anomalías; por ejemplo, la separación de la información general de empleados de
los datos específicos de médicos, y el uso de tablas puente para relaciones
muchos-a-muchos.

\begin{figure}
    \caption{Ejemplo de normalización aplicada}
    \centering
    % \includegraphics[width=\linewidth]{img/normalizacion.png}
    \fbox{\parbox[c][5cm][c]{\linewidth}{\centering [Captura de pantalla: antes y después de la normalización]}}
    \label{fig:norm}
\end{figure}

% =====================================================================
\section{Implementación en MySQL}
% =====================================================================
Las tablas se crearon mediante un script DDL y se poblaron con datos de prueba
mediante un script DML. La siguiente figura muestra las tablas creadas en el
manejador.

\begin{figure}
    \caption{Tablas creadas en MySQL}
    \centering
    % \includegraphics[width=\linewidth]{img/tablas_mysql.png}
    \fbox{\parbox[c][5cm][c]{\linewidth}{\centering [Captura de pantalla: tablas en MySQL]}}
    \label{fig:tablas}
\end{figure}

\begin{figure}
    \caption{Datos insertados en una tabla}
    \centering
    % \includegraphics[width=\linewidth]{img/datos_mysql.png}
    \fbox{\parbox[c][5cm][c]{\linewidth}{\centering [Captura de pantalla: registros insertados]}}
    \label{fig:datos}
\end{figure}

% =====================================================================
\section{Vistas de Reporte}
% =====================================================================
Se implementaron las vistas requeridas: historial clínico completo, inventario
con alertas de stock mínimo y medicamentos más recetados y surtidos. La
siguiente figura muestra el resultado de la Vista 3.

\begin{figure}
    \caption{Resultado de la Vista 3: medicamentos más recetados y surtidos}
    \centering
    % \includegraphics[width=\linewidth]{img/vista3.png}
    \fbox{\parbox[c][5cm][c]{\linewidth}{\centering [Captura de pantalla: Vista 3]}}
    \label{fig:vista3}
\end{figure}

% =====================================================================
\section{Conexión con la Aplicación}
% =====================================================================
Se desarrolló una aplicación Java que se conecta a la base de datos mediante
JDBC y consulta la información. La figura muestra la ejecución exitosa.

\begin{figure}
    \caption{Ejecución de la aplicación Java (JDBC)}
    \centering
    % \includegraphics[width=\linewidth]{img/app_java.png}
    \fbox{\parbox[c][5cm][c]{\linewidth}{\centering [Captura de pantalla: salida de App.java]}}
    \label{fig:app}
\end{figure}

% =====================================================================
\section{Conclusión}
% =====================================================================
El diseño e implementación de la base de datos cumple con los requerimientos del
caso de negocio, permitiendo un control integrado de consultas, recetas,
inventario y ventas, y sentando las bases para reportes de valor para la gestión
del consultorio y la farmacia.

\end{document}
